%% LyX 2.4.4 created this file.  For more info, see https://www.lyx.org/.
%% Do not edit unless you really know what you are doing.
\documentclass[english]{article}
\usepackage[T1]{fontenc}
\usepackage[latin9]{inputenc}
\usepackage{enumitem}
\usepackage{amsmath}
\usepackage{amsthm}
\usepackage{amssymb}
\usepackage{cancel}
\usepackage{geometry}
\geometry{verbose,tmargin=1in,bmargin=1in,lmargin=1in,rmargin=1in}
\PassOptionsToPackage{normalem}{ulem}
\usepackage{ulem}

\makeatletter
%%%%%%%%%%%%%%%%%%%%%%%%%%%%%% Textclass specific LaTeX commands.
\numberwithin{equation}{section}
\numberwithin{figure}{section}
\newlength{\lyxlabelwidth}      % auxiliary length 

%%%%%%%%%%%%%%%%%%%%%%%%%%%%%% User specified LaTeX commands.
\usepackage{fancyhdr}
\usepackage{lastpage}
\pagestyle{fancy}
\usepackage{enumitem}
\usepackage{ifthen}
\everymath=\expandafter{\the\everymath\displaystyle}
%\DeclareMathSizes{10}{10}{10}{10}
\usepackage{multicol}

\usepackage{tikz}
\usetikzlibrary{patterns}
\usetikzlibrary{plotmarks}
\usepackage{pgfplots}
\pgfplotsset{compat=newest}

%\usepackage{advdate}    % Advancing/saving dates
\usepackage{datetime}   % Dates formatting
%\usdate
%\usepackage{datenumber} % Counters for dates


%\usepackage{calc}
%\usepackage[nomessages]{fp}
% Added by lyx2lyx
\setlength{\parskip}{\smallskipamount}
\setlength{\parindent}{0pt}

\makeatother

\usepackage{babel}
\begin{document}
\renewcommand{\author}{Dr.\,James Ashe}

\newcommand{\classNum}{335}

\newcommand{\classSec}{A}

\newcommand{\examType}{Final Exam}

\renewcommand{\date}{\formatdate{17}{5}{2013}}

%%First page's header%%%%%%%%%%%%%%%%%%%%%%
\fancypagestyle{firststyle} %%header settings for first pagr
{
	\fancyhead[L]{MTH \classNum{}(\classSec{}) \author{}\\
	\textbf{\examType{}} \date{}}
	\fancyhead[C]{\textbf{Solutions}}
	\setlength{\headsep}{14pt}
}
%%%%%%%%%%%%%%%%%%%%%%%%%%%%%%%%%%%%%%%%%%

%%Next page's header%%%%%%%%%%%%%%%%%%%%%%%
\setlist{nolistsep}
\lhead{\classNum{}(\classSec{})} 
\chead{\textbf{}} 
\cfoot{\thepage{}~of~\pageref{LastPage}} 
\setlength{\headsep}{5pt} 
%%%%%%%%%%%%%%%%%%%%%%%%%%%%%%%%%%%%%%%%%%%

%%Various settings; see the preamble for other settings%%%%%
%\setenumerate[0]{leftmargin=12pt,itemindent=0pt,labelwidth=0pt}
\setlist{topsep=.5pt}
\renewcommand{\labelenumi}{\textbf{\arabic{enumi}.}}
\renewcommand{\labelenumii}{\textbf{\alph{enumii}.}}
\thispagestyle{firststyle}
%%%%%%%%%%%%%%%%%%%%%%%%%%%%%%%%%%%%%%%%%%

%\newcommand{\circled}[1]{\raisebox{.5pt}{\textcircled{\raisebox{-.9pt} {#1}}}} % a simple circle command

\newcommand*\circled[1]{\tikz[baseline=(char.base)]{     \node[shape=circle,draw,inner sep=-.75pt] (char) {#1};}} %change sep for larger circle
\newcommand{\ndiv}{\hspace{-4pt}\not|\hspace{2pt}} 
\newcommand{\dspace}{\vspace{.75cm}} 

\footnotesize
\begin{enumerate}
\item Let $A=\left\{ 1,2,3,5\right\} $ and $B=\left\{ a,b,2,c,5\right\} $

\begin{enumerate}
\item Find $A\cap B$

\begin{enumerate}
\item []$=\left\{ 2,5\right\} $
\end{enumerate}
\item Find $A\cup B$.

\begin{enumerate}
\item []$=\left\{ 1,2,3,5,a,b,c\right\} $
\end{enumerate}
\item Find $B-A$.

\begin{enumerate}
\item []$=\left\{ a,b,c\right\} $
\end{enumerate}
\item Find $A-B$.

\begin{enumerate}
\item []$=\left\{ 1,3\right\} $
\end{enumerate}
\end{enumerate}
\item Let $A$ be a nonempty subset of $B$. Which of the following is true?

\begin{enumerate}
\item [\cancel{\textbf{(a)}}]$A\cap B=B$.
\item [\cancel{\textbf{(b)}}]$A\cup B=A$.
\item [\circled{\textbf{(c)}}]$A\cup B=B$.
\item [\cancel{\textbf{(d)}}]$A\cap B=\emptyset$.
\end{enumerate}
\item Let $A=\left\{ 3,5\right\} $ and $B=\left\{ a,b,5\right\} $. 

\begin{enumerate}
\item List all possible subsets of $B$.

\begin{enumerate}
\item []$\left\{ \emptyset\right\} ,\left\{ a\right\} ,\left\{ b\right\} ,\left\{ c\right\} ,\left\{ a,b\right\} ,\left\{ a,5\right\} ,\left\{ b,5\right\} ,\left\{ a,b,5\right\} $
\end{enumerate}
\item Find $A\times B$. 

\begin{enumerate}
\item []$\left\{ (3,a),(3,b),(3,5),(5,a),(5,b),(5,5)\right\} $
\end{enumerate}
\end{enumerate}
\item Prove that $A\cup\emptyset=A$.

\begin{proof}
$\emptyset\subseteq A$ and $A\subseteq A$ thus $A\cup\emptyset\subseteq A$
and $A\subseteq A\cup\emptyset$.

Alternatively, 

$\left(A\subseteq A\cup\emptyset\right)$ Let $a\in A$ then $a\in A$
or $\emptyset$ thus $a\in A\cup\emptyset$.

$\left(A\cup\emptyset\subseteq A\right)$ If $A=\emptyset$ then $A\cup\emptyset=\emptyset\subseteq A$.
Otherwise if $a\in A\cup\left\{ \emptyset\right\} $ then $a\in A$.
Thus $A\cup\left\{ \emptyset\right\} =A$.
\end{proof}
\item Consider the map $f:\mathbb{R}\times\mathbb{R}\rightarrow\mathbb{R}$
where $f(x,y)=xy$. Prove that $f$ is onto.

\begin{proof}
Let $z\in\mathbb{R}$. $1\in\mathbb{R}$ so then $f(1,z)=1\cdot z=z$.
Thus $f$ is onto.
\end{proof}
\item Define the relation $\mathcal{R}$ on the real number set $\mathbb{R}$
as $x\mathcal{R}y\iff xy\geq0$. 

\begin{enumerate}
\item Show that $\mathcal{R}$ is reflexive and symmetric.

\begin{proof}
Reflexive: $x\cdot x=x^{2}\geq0$. 

Symmetric: If $x\cdot y\geq0$ then $x\cdot y=y\cdot x\geq0$.
\end{proof}
\item Is $\mathcal{R}$ an equivalence relation? Explain.

\begin{enumerate}
\item []No. $\mathcal{R}$ is not transitive. $-1\cdot0\geq0$ and $0\cdot1\geq0$
but $-1\cdot1=-1<0$.
\end{enumerate}
\end{enumerate}
\item Let $X$ be a nonempty set. Define the following relation on the set
of subsets of $X$ as follows $A\mathcal{R}B\iff A\cap B\neq\emptyset$
for $A,B\subseteq X$. 

\begin{enumerate}
\item Show that $\mathcal{R}$ is symmetric.

\begin{proof}
Suppose that $A\cap B\neq\emptyset$ then $B\cap A\neq\emptyset$.
\end{proof}
\item Show that $\mathcal{R}$ is not reflexive.

\begin{proof}
$\emptyset\in X$, and $\emptyset\cap\emptyset=\emptyset$.
\end{proof}
\end{enumerate}
\item Let $f:A\rightarrow B$ and $g:B\rightarrow C$ be two functions.

\begin{enumerate}
\item If $g\circ f$ is one-to-one then 

\begin{enumerate}
\item [\cancel{\textbf{(a)}}]$f$ is onto.
\item [\textbf{(b)}]$f$ is one-to-one.
\item [\textbf{(c)}]$g$ is one-to-one.
\item [\circled{\textbf{(d)}}]Both $f$ and $g$ are one-to-one.
\item [\cancel{\textbf{(e)}}]None of the above. 
\end{enumerate}
\item If $\left|A\right|>\left|B\right|$ and $A\neq\emptyset$ can $f$
be one-to-one?

\begin{enumerate}
\item []No. By the Pigeon Hole Theorem.
\end{enumerate}
\end{enumerate}
\item Let $S_{3}$ be the symmetric group on three elements. 

\begin{enumerate}
\item Which of the following is true?

\begin{enumerate}
\item [\textbf{(a)}]$S_{3}$ has at least one non-abelain subgroup.
\item [\cancel{\textbf{(b)}}]All proper subgroups of $S_{3}$ are abelian.
\item [\cancel{\textbf{(c)}}]$S_{3}$ is abelian.
\item [\cancel{\textbf{(d)}}]$S_{3}$ is cyclic.
\end{enumerate}
\item What is the order of $S_{3}$ and $S_{4}$?

\begin{enumerate}
\item []$\left|S_{3}\right|=3!=6$
\item []$\left|S_{4}\right|=4!=24$
\end{enumerate}
\end{enumerate}
\item Let $f:G_{1}\rightarrow G_{2}$ be a group homomorphism and let $H\leq G_{2}$.
Recall that $f^{-1}(H)=\left\{ g\in G_{1}|f(g)\in H\right\} $. 

\begin{enumerate}
\item Which of the following is true?

\begin{enumerate}
\item [\cancel{\textbf{(a)}}]$f^{-1}(H)$ consists only of the identity
of $G_{1}$.
\item [\circled{\textbf{(b)}}]$f^{-1}(H)$ \uline{is} a subgroup of $G_{1}$.
\item [\cancel{\textbf{(c)}}]$gf^{-1}(H)=f^{-1}(H)g$ for all $g\in G_{1}$
(the left and right cosets are equal).
\item [\cancel{\textbf{(d)}}]$f^{-1}(H)$ \uline{is not} a subgroup of
$G_{1}$
\end{enumerate}
\item Let $a,b\in G_{1}$. Does $f(ab)=f(a)f(b)$? 

\begin{enumerate}
\item []Yes. Since $f$ is homomorphic $f(ab)=f(a)f(b)$ by definition.
\end{enumerate}
\end{enumerate}
\item Let $G$ be a finite group of order $12$ and $g\in G$. 

\begin{enumerate}
\item The order of $g$ \uline{cannot }be:

\begin{enumerate}
\item [\circled{\textbf{(a)}}]8
\item [\cancel{\textbf{(b)}}]6
\item [\cancel{\textbf{(c)}}]4
\item [\cancel{\textbf{(d)}}]3
\end{enumerate}
\item Let $\left|H\right|=6$. Does Lagrange's theorem imply that $H\leq G$?

\begin{enumerate}
\item []No. Lagrange's Theorem is a necessary but not sufficient condition
for a subset to be a subgroup. For example, the alterating group $A_{4}$
has $12$ elements but no subgroup of order $6$.
\end{enumerate}
\end{enumerate}
\item Let $G$ be a group of order $4$. Which of the following is true?

\begin{enumerate}
\item [\cancel{\textbf{(a)}}]$G$ is a cyclic group.
\item [\cancel{\textbf{(b)}}]$G$ is $\mathbb{Z}_{2}\times\mathbb{Z}_{2}$.
\item [\cancel{\textbf{(c)}}]$G$ is $\mathbb{Z}_{4}$.
\item [\circled{\textbf{(d)}}]$G$ is abelian.
\end{enumerate}
\item Let $G=\left\{ e,a_{1},a_{2},a_{3},a_{4},a_{5}\right\} $ be a group
where $e$ is the identity.

\begin{enumerate}
\item Find a subgroup of $G$.

\begin{enumerate}
\item []The trivial group $e$ is a subgroup of every group $G$.
\end{enumerate}
\item Find a subgroup of $G$ distinct from the group given for part \textbf{(a)}.

\begin{enumerate}
\item []$\left\langle a_{i}\right\rangle \leq G$, and $\left\langle a_{i}\right\rangle \neq e$
for any $1\leq i\leq5$.
\end{enumerate}
\item Can $\left|a_{2}\right|=5$?

\begin{enumerate}
\item []No. By Lagrange's theorem $\left|a_{2}\right|\bigr|\left|G\right|$
and $5\ndiv6$
\end{enumerate}
\end{enumerate}
\item Let $H$ and $K$ be subgroups be non-trivial subgroups of $G$. Prove
or disprove that $H\cap K$ is a subgroup of $G$.

\begin{proof}
$H\cap K\leq G$. Let $a,b\in H\cap K$. Then $b^{-1}\in H$ and $K$
since $b\in H$ and $K$ and $H,K\leq G$. Thus $ab^{-1}\in H\cap K$
because $a\in H$ and $K$ and subgroups are closed under their operation.
\end{proof}
\item Let $R$ be a ring and $a,b\in R$. 

\begin{enumerate}
\item Which statement is true?

\begin{enumerate}
\item [\cancel{\textbf{(a)}}]If $ab=0$ then $a=0$ or $b=0$.
\item [\cancel{\textbf{(b)}}]If $a^{2}=0$ then $a=0$.
\item [\cancel{\textbf{(c)}}]If $a\neq0$ and $b\neq0$ then $ab\neq0$.
\item [\circled{\textbf{(d)}}]None of the above.
\end{enumerate}
\item Show that $\mathbb{Q}$, the set of rational numbers, is a ring under
addition and multiplication. 

\begin{proof}
$\mathbb{Q}$ is an abelian group. $\mathbb{Q}\subset\mathbb{R}$
so then associativity and distribution hold under multiplication.
Thus $\mathbb{Q}$ is a ring.
\end{proof}
\end{enumerate}
\end{enumerate}

\end{document}
